\chapter{Introdução}
\label{cap:introducao}

\section{Contexto e Problema}

O setor da construção civil no Brasil figura como um dos mais críticos em relação à segurança laboral, ocupando a 6ª posição no ranking de acidentes de trabalho típicos e a 9ª em casos de doenças profissionais ou ocupacionais, segundo dados do \gls{MTE} \cite{mte2024esocial}. Em 2023, o país registrou 2.888 mortes decorrentes de acidentes laborais e 499.955 ocorrências de acidentes típicos via sistema eSocial, com a construção liderando esse triste cenário. Em 2024, os números agravaram-se: foram 724.228 acidentes de trabalho e 3.394 óbitos, um aumento de 11,16\% em relação ao ano anterior, com a ocupação de pedreiro entre as três com maior número de mortes \cite{mte2025canpat}. Entre as principais causas estão impactos com objetos, quedas de altura, choques elétricos e soterramentos ou desmoronamentos, fatores que frequentemente estão associados à falta de equipamentos de proteção adequados, falhas na fiscalização e deficiências na formação técnica.

Restringir o recorte ao subsetor de construção de edifícios (\gls{CNAE} 4120) torna o quadro mais nítido. Esse subsetor concentra 39,6\% dos afastamentos acidentários da construção (benefício por incapacidade B91), sendo o segmento mais afetado, segundo o Observatório de Segurança e Saúde no Trabalho \cite{smartlab2024}. O perfil de lesão é dominado por fraturas, que respondem por 51,8\% dos afastamentos por capítulo da Classificação Internacional de Doenças e por 67\% das lesões por causas externas — padrão coerente com eventos de alta energia, como quedas de altura e soterramentos \cite{smartlab2024,moreira2024accidents}. As ocupações mais atingidas são servente de obras (20\%) e pedreiro (16,3\%), exatamente o público-alvo dos treinamentos de segurança em canteiro.

No âmbito regional, os dados do Ceará reforçam a urgência do problema. Em 2024, a construção civil cearense registrou 593 acidentes de trabalho e 12 óbitos, com a fratura como lesão mais frequente (135 casos) e o subsetor de construção de edifícios à frente dos registros \cite{sesifiec2024}. Enquanto os acidentes na construção recuaram no país, no Ceará houve aumento de 2,87\% em relação ao ano anterior, na contramão da tendência nacional. Servente de obras (118 casos) e pedreiro (53 casos) lideram as ocupações atingidas, confirmando localmente o perfil profissional mais exposto.

A \gls{NR}-01 \cite{brasil2024nr01}, em seu tópico 1.7, destaca a importância da capacitação contínua em \gls{SST}, exigindo atualizações periódicas para preparar os trabalhadores para os riscos ocupacionais, uma abordagem reforçada pela NR-18 \cite{brasil2025nr18}, que foca na prevenção de acidentes e na criação de ambientes seguros na construção civil, e pela NR-35 \cite{brasil2023nr35}, específica para o trabalho em altura — risco central no perfil de acidentes aqui descrito. Para atender a essas exigências, o treinamento em \gls{RV} surge como uma alternativa tecnológica que permite aos trabalhadores praticarem, de forma imersiva e segura, a identificação de riscos, o uso correto de \glspl{EPI} e \glspl{EPC} e a resposta a emergências, sem exposição a perigos reais. Essa tecnologia não só melhora a retenção do conhecimento, mas também fortalece a cultura de segurança, contribuindo para a redução de acidentes e o cumprimento das normas regulamentadoras.

\section{Objetivo Geral}

Desenvolver e validar a Fase 1 de um protótipo modular de Realidade Virtual Imersiva para treinamento de segurança em canteiros de obras, focando na identificação e mitigação de riscos relacionados a trabalhos em altura e uso de Equipamentos de Proteção Individual (EPIs), no contexto brasileiro.

\section{Hipótese/Pressuposto}

\begin{itemize}
    \item \textbf{H1 -- Eficácia do treinamento:} O treinamento em RV resulta em ganho significativo de conhecimento sobre procedimentos de segurança em trabalhos em altura e uso correto de EPIs, evidenciado pela comparação entre pré e pós-testes nos grupos experimental e controle.

    \item \textbf{H2 -- Motivação e engajamento:} O treinamento em RV promove maior motivação e engajamento dos participantes em comparação com métodos tradicionais, evidenciado pela comparação dos questionários de motivação entre os dois grupos.

    \item \textbf{H3 -- Qualidade da experiência imersiva:} O sistema de RV apresenta níveis adequados de usabilidade, conforme medido pelo \gls{SUS}, e gera sensação significativa de presença no ambiente virtual, conforme avaliado pelo \gls{IPQ}, junto ao grupo experimental.
\end{itemize}
